%======================================================================
% PRESENTACIÓN - SSSIFANB Manual de Usuario
%======================================================================
% Introducción al manual y al sistema
%======================================================================

\chapter*{Presentación}
\addcontentsline{toc}{chapter}{Presentación}

\section*{Acerca de este Manual}

El presente documento constituye la guía oficial de usuario para el \textbf{Sistema de Seguridad Social de la Fuerza Armada Nacional Bolivariana (SSSIFANB)}, desarrollado bajo el marco del \textbf{Instituto de Previsión Social de la Fuerza Armada Nacional Bolivariana (IPSFANB)}.

Este manual ha sido elaborado siguiendo los lineamientos establecidos en la normativa legal vigente, particularmente la \textbf{Ley del Sistema de Seguridad Social de la Fuerza Armada Nacional Bolivariana}, y tiene como propósito fundamental facilitar el proceso de aprendizaje y dominio de las funcionalidades que ofrece el sistema a todos los usuarios autorizados.

\section*{Objetivos del Manual}

\begin{itemize}[leftmargin=*]
    \item \textbf{Proporcionar}: Una guía completa de las funcionalidades del sistema
    \item \textbf{Facilitar}: El proceso de aprendizaje para nuevos usuarios
    \item \textbf{Optimizar}: El tiempo de consulta mediante una estructura intuitiva
    \item \textbf{Estandarizar}: Los procedimientos operativos del sistema
\end{itemize}

\section*{Audiencia Destinataria}

Este manual está dirigido a:

\begin{itemize}[leftmargin=*]
    \item \textbf{Afiliados}: Miembros de la Fuerza Armada en situación de actividad o retiro
    \item \textbf{Beneficiarios}: Familiares y dependientes legales de los afiliados
    \item \textbf{Administradores}: Personal del IPSFANB responsable de la gestión del sistema
    \item \textbf{Operadores}: Funcionarios autorizados para procesar solicitudes
\end{itemize}

\section*{Estructura del Documento}

El manual está organizado en las siguientes secciones principales:

\begin{enumerate}[leftmargin=*]
    \item \textbf{Fundamentos y Primeros Pasos}: Conceptos básicos y proceso de autenticación
    \item \textbf{Módulos del Sistema}: Descripción detallada de cada módulo funcional
    \item \textbf{Administración}: Configuración y gestión del sistema
    \item \textbf{Anexos}: Referencias técnicas adicionales
\end{enumerate}

\section*{Convenciones Tipográficas}

\begin{itemize}[leftmargin=*]
    \item \textbf{\texttt{Texto en código}}: Representa nombres de campos, botones o rutas del sistema
    \item \textbf{\textit{Texto en itálica}}: Indica términos o conceptos importantes
    \item \textbf{\textbf{Negrita}}: Destaca instrucciones o advertencias críticas
\end{itemize}

\begin{figure}[H]
\centering
\begin{tcolorbox}[
    enhanced,
    colback=sandraSoft,
    colframe=sandraGreen,
    width=0.9\textwidth,
    halign=center
]
\small
\textbf{Nota Importante}\\
\smallskip
Para garantizar la seguridad de la información, el sistema requiere autenticación de dos factores (2FA) en todos los accesos. Asegúrese de tener sus credenciales y dispositivo autenticador disponibles antes de iniciar sesión.
\end{tcolorbox}
\end{figure}

\vfill

\begin{flushright}
\textbf{Fecha de publicación:} Marzo 2026\\
\textbf{Versión del documento:} 1.0.0
\end{flushright}

\newpage
